\chapter{系统设计}

\section{总体架构}
FTCL 的总体架构如图 \ref{fig:system-arch} 所示。一个脚本首先被解析为 commands 和 words，随后由 interpreter 对 word 进行求值，完成变量替换、命令替换、数组访问和字符串拼接，最后将第一个求值结果作为 command name 进行分发。\texttt{geom} command 位于语言层和计算层之间：它负责解析几何参数或 UVec handle，并根据用户指定的 device 选择 CPU backend 或 CUDA backend。UVec 则处在 command layer 和 backend layer 之间，负责维护跨设备数据的位置、有效性和同步过程。

\begin{figure}[H]
  \centering
  \includegraphics[width=0.96\textwidth]{system-architecture.png}
  \caption{FTCL 几何脚本系统总体架构（FTCL geometry scripting system architecture）}
  \label{fig:system-arch}
\end{figure}

这个设计的核心是分层而不是堆叠。Parser 和 interpreter 专注于语言语义；command layer 提供稳定的用户接口；UVec 负责 data residency 和 synchronization；CPU/CUDA backend 只实现具体算法。这样做的好处是：语言可以继续扩展，几何算法可以继续增加，CUDA 路径也可以在不破坏脚本接口的情况下逐步优化。

\section{工程规模}
表 \ref{tab:engineering-scope} 总结了项目的主要工程组成。列出这张表的目的不是简单展示文件数量，而是让系统边界变得清楚：FTCL 不是几个孤立的几何函数，而是由 interpreter、memory abstraction、geometry library、CUDA backend、testing framework 和 benchmark pipeline 共同组成的完整原型系统。

\begin{table}[htbp]
\centering
\caption{FTCL 系统工程组成（Engineering scope）}
\label{tab:engineering-scope}
\small
\begin{tabularx}{\textwidth}{lXXX}
\toprule
Component & Main files & Main classes/functions & Role \\
\midrule
Interpreter core & \texttt{interp.hpp}, \texttt{type.hpp}, \texttt{parser.hpp} & \texttt{Interp}, \texttt{Value}, \texttt{Word}, \texttt{Script} & 负责 script parsing、evaluation、variables、procedures 和 command dispatch。 \\
Expression parser & \texttt{expr.hpp} & \texttt{ExprParser}, \texttt{Datum} & 支持 arithmetic、comparison、logical、bitwise、ternary 以及 substitution-aware expressions。 \\
UVec layer & \texttt{uvec.hpp}, \texttt{commands.hpp} & \texttt{UVec<T>}, pointer wrappers, handle manager & 提供 cross-device vector storage、CPU/CUDA synchronization 和 script-visible handles。 \\
Geometry CPU backend & \texttt{geometry.hpp} & \texttt{Point}, \texttt{Segment}, \texttt{BoundingBox}, batch algorithms & 在 CPU data 上实现 scalar 和 batch two-dimensional geometry algorithms。 \\
CUDA backend & \texttt{geometry\_cuda.*} & kernels and launch wrappers & 为 independent-work-item geometry algorithms 提供 GPU acceleration。 \\
Tests & \texttt{test/*.cpp} & semantic, equivalence, random, edge, stress tests & 验证 language semantics、UVec consistency 和 CPU/GPU correctness。 \\
Benchmarks and figures & \texttt{docs/data}, \texttt{docs/scripts}, \texttt{docs/figures} & CSV files and Python generators & 生成 reproducible performance data 和 thesis-ready figures。 \\
\bottomrule
\end{tabularx}
\end{table}

\begin{table}[htbp]
\centering
\caption{实现规模统计（Quantitative implementation statistics）}
\label{tab:implementation-statistics}
\small
\begin{tabularx}{\textwidth}{lrX}
\toprule
Metric & Value & Meaning \\
\midrule
C++/CUDA source files & 56 & \texttt{include}、\texttt{src} 和 \texttt{test} 下的 header、source、CUDA、benchmark 和 test files。 \\
C++/CUDA source lines & 19151 & \texttt{.hpp}、\texttt{.cpp} 和 \texttt{.cu} 文件总行数。 \\
Header/interface lines & 12881 & 核心 interpreter、UVec、parser、expression、command 和 geometry interfaces。 \\
CUDA/source implementation lines & 642 & CUDA backend 和非 header source files。 \\
Test source lines & 5628 & Regression、equivalence、edge-case、random、stress 和 benchmark-related test code。 \\
CTest entries & 31 & Build-system registered tests，用于 regression 和 validation。 \\
CUDA kernels & 17 & CUDA geometry backend 以及相关 test path 中的 device kernels。 \\
Benchmark executables & 6 & 用于 runtime、geometry 和 multi-GPU experiments 的 non-interactive benchmark programs。 \\
Benchmark/figure scripts & 5 & \texttt{docs/scripts} 下的数据和图表生成脚本。 \\
Generated SVG figures & 18 & \texttt{docs/figures} 下的 reproducible vector figures。 \\
Thesis PNG figures & 18 & 复制到 \texttt{contents/figures} 中用于论文编译的 rendered figures。 \\
\bottomrule
\end{tabularx}
\end{table}

\section{Interpreter Runtime}
Runtime 由 values、words、scripts、scopes 和 commands 组成。Value 表示脚本层的值；Word 表示一个 command argument 在求值前后的结构；Script 是 command list；Interpreter 维护 command tables 和 variable scopes。FTCL 的执行路径保持 Tcl-style 的简单性，但在内部为后续扩展保留了结构化表示。

一个普通脚本的执行过程如下：

\begin{enumerate}
  \item 将 source text 解析为 script structure。
  \item 按顺序执行一个个 command。
  \item 对每个 word 执行 variable substitution、command substitution、array lookup 和 string concatenation。
  \item 将第一个求值后的 word 作为 command name。
  \item 调用对应的 native C++ command function，并返回 value 或 exception。
\end{enumerate}

这种结构有一个重要优势：新能力不需要发明新语法。\texttt{thread}、\texttt{uvec} 和 \texttt{geom} 都是 command，而不是语言核心的特殊分支。因此 FTCL 可以保持小而清晰，同时承载并发、跨设备内存和几何算法。

\section{词法分析流程（Lexical Analysis Flow）}
Lexical layer 将原始 script string 转换为 token stream。FTCL 的 token-stream 方向会保留 token text 和 token boundaries，使后续阶段能明确识别 command endings、comments、whitespace、braces、brackets、variable references 和 ordinary text。算法 \ref{alg:lexer} 描述了 lexer 的基本流程。

\begin{algorithm}[H]
\caption{FTCL Lexical Analysis}
\label{alg:lexer}
\KwInput{source string $source$}
\KwOutput{token stream $tokens$}
$tokens \leftarrow [\ ]$\;
$cursor \leftarrow 0$\;
\While{$cursor < |source|$}{
  $ch \leftarrow source[cursor]$\;
  \uIf{$ch \in \{\mathrm{space},\mathrm{tab}\}$}{
    $start \leftarrow cursor$\;
    consume horizontal whitespace\;
    append \textsc{Token}(\textsc{Whitespace}, $source[start:cursor]$) to $tokens$\;
  }
  \uElseIf{$ch$ is a newline or $ch = \texttt{;}$}{
    append \textsc{Token}(\textsc{CommandEnd}, $ch$) to $tokens$\;
    $cursor \leftarrow cursor + 1$\;
  }
  \uElseIf{$ch = \texttt{\#}$ and $cursor$ is at command start}{
    $start \leftarrow cursor$\;
    consume until newline or end of file\;
    append \textsc{Token}(\textsc{Comment}, $source[start:cursor]$) to $tokens$\;
  }
  \uElseIf{$ch$ is a Tcl special delimiter}{
    append \textsc{Token}(\textsc{Special}, $ch$) to $tokens$\;
    $cursor \leftarrow cursor + 1$\;
  }
  \Else{
    $start \leftarrow cursor$\;
    consume ordinary characters until a token boundary\;
    append \textsc{Token}(\textsc{Text}, $source[start:cursor]$) to $tokens$\;
  }
}
append \textsc{Token}(\textsc{EndOfFile}, empty string) to $tokens$\;
\KwReturn{$tokens$}\;
\end{algorithm}

Lexer 不执行 substitution。它只记录输入中可观察的结构。这样做非常重要，因为 Tcl-style substitution 本质上是 semantic action：\texttt{\$} 可能需要 variable lookup，\texttt{[...]} 可能需要 recursive command evaluation。把 tokenization 和 evaluation 分开，可以让 parser 更稳定，也便于错误定位和回归测试。

\section{语法分析流程（Syntax Analysis Flow）}
Parser 消费 token stream，并把 tokens 分组成 commands 和 words。Command 在 newline、semicolon 或 end-of-file 处结束，但如果 parser 位于 braces、quotes 或 brackets 内部，则这些符号不再表示 command boundary。Parser 不执行 command，它只构造后续可执行的 script representation。

\begin{algorithm}[H]
\caption{Token-stream Command Parsing}
\label{alg:parse-script}
\KwInput{token stream $tokens$}
\KwOutput{script object $script$}
\Proc{\textsc{ParseScript}($tokens$)}{
  $cursor \leftarrow \textsc{TokenCursor}(tokens)$\;
  $script \leftarrow [\ ]$\;
  \While{not $cursor.atEnd()$}{
    $cursor.skipWhitespaceAndComments()$\;
    \If{$cursor.atCommandEnd()$}{
      $cursor.next()$\;
      \textbf{continue}\;
    }
    $command \leftarrow \textsc{ParseCommand}(cursor)$\;
    append $command$ to $script.commands$\;
    \uIf{$cursor.atCommandEnd()$}{
      $cursor.next()$\;
    }
    \uElseIf{not $cursor.atEnd()$}{
      raise ``expected command boundary''\;
    }
  }
  \KwReturn{$script$}\;
}
\Proc{\textsc{ParseCommand}($cursor$)}{
  $words \leftarrow [\ ]$\;
  \While{not $cursor.atEnd()$ and not $cursor.atCommandEnd()$}{
    $cursor.skipInlineWhitespace()$\;
    \If{$cursor.atCommandEnd()$ or $cursor.atEnd()$}{
      \textbf{break}\;
    }
    $word \leftarrow \textsc{ParseWord}(cursor)$\;
    append $word$ to $words$\;
  }
  \KwReturn{\textsc{Command}($words$)}\;
}
\end{algorithm}

Word parser 比 command parser 更细，因为一个 word 内部可能包含 literal text、variable references、command substitutions、array references、braced literals、quoted strings 和 expansion syntax。关键设计是：parser 将这些组成部分保存为结构化 word nodes，而 interpreter 负责决定它们在运行时的 value。

\begin{algorithm}[H]
\caption{Word Parsing}
\label{alg:parse-word}
\KwInput{token cursor $cursor$}
\KwOutput{word node}
\Proc{\textsc{ParseWord}($cursor$)}{
  \uIf{$cursor$ matches an open brace}{
    \KwReturn{\textsc{Word.Value}(\textsc{ReadBalancedBracedText}($cursor$))}\;
  }
  \uElseIf{$cursor$ matches a quote}{
    $parts \leftarrow [\ ]$\;
    \While{$cursor$ does not match the closing quote}{
      append \textsc{ParseWordPart}($cursor$, true) to $parts$\;
    }
    \KwReturn{\textsc{Word.Tokens}($parts$)}\;
  }
  \uElseIf{$cursor$ matches the expansion prefix \texttt{\{*\}}}{
    $child \leftarrow \textsc{ParseWord}(cursor)$\;
    \KwReturn{\textsc{Word.Expand}($child$)}\;
  }
  $parts \leftarrow [\ ]$\;
  \While{not $cursor.atWordBoundary()$}{
    append \textsc{ParseWordPart}($cursor$, false) to $parts$\;
  }
  \uIf{$parts$ contains one literal part}{
    \KwReturn{\textsc{Word.Value}($parts[0]$)}\;
  }
  \Else{
    \KwReturn{\textsc{Word.Tokens}($parts$)}\;
  }
}
\Proc{\textsc{ParseWordPart}($cursor$, $quoted$)}{
  \uIf{$cursor.peek() = \texttt{\$}$}{
    \KwReturn{\textsc{ParseVariableReference}($cursor$)}\;
  }
  \uElseIf{$cursor.peek() = \texttt{[}$}{
    \KwReturn{\textsc{Word.Script}(\textsc{ReadBalancedBracketScript}($cursor$))}\;
  }
  \uElseIf{$cursor.peek()$ is a backslash}{
    \KwReturn{\textsc{Word.String}(\textsc{ParseBackslashSequence}($cursor$))}\;
  }
  \Else{
    \KwReturn{\textsc{Word.String}(\textsc{ReadLiteralText}($cursor$, $quoted$))}\;
  }
}
\end{algorithm}

\section{语义执行流程（Semantic Execution Flow）}
Semantic execution 发生在 parsing 之后。Interpreter 逐条执行 command，对每个 word 求值，并把第一个 value 作为 command name，剩余 values 作为 arguments。Command function 才真正改变变量、创建 UVec、启动 CUDA kernel 或返回结果。

\begin{algorithm}[H]
\caption{Command-level Semantic Execution}
\label{alg:semantic-exec}
\KwInput{interpreter $interp$, parsed script $script$}
\KwOutput{last command result}
\Proc{\textsc{EvalScript}($interp$, $script$)}{
  $result \leftarrow$ empty string\;
  \ForEach{$command \in script.commands$}{
    \If{$command.words$ is empty}{
      \textbf{continue}\;
    }
    $values \leftarrow [\ ]$\;
    \ForEach{$word \in command.words$}{
      $value \leftarrow \textsc{EvalWord}(interp, word)$\;
      \uIf{$word$ is an expansion word}{
        extend $values$ by \textsc{ParseList}($value$)\;
      }
      \Else{
        append $value$ to $values$\;
      }
    }
    $name \leftarrow values[0].asString()$\;
    $proc \leftarrow interp.lookupCommand(name)$\;
    \If{$proc$ is missing}{
      raise ``invalid command name''\;
    }
    $result \leftarrow proc(interp, values)$\;
    \If{$result.code$ is \textsc{Return}, \textsc{Break}, \textsc{Continue}, or \textsc{Error}}{
      propagate $result$ according to the current context\;
    }
  }
  \KwReturn{$result$}\;
}
\end{algorithm}

\begin{algorithm}[H]
\caption{Word Evaluation}
\label{alg:word-eval}
\KwInput{interpreter $interp$, word node $word$}
\KwOutput{runtime value}
\Proc{\textsc{EvalWord}($interp$, $word$)}{
  \uIf{$word.kind = \textsc{Value}$}{
    \KwReturn{$word.literal$}\;
  }
  \uElseIf{$word.kind = \textsc{String}$}{
    \KwReturn{$word.text$}\;
  }
  \uElseIf{$word.kind = \textsc{VarRef}$}{
    \KwReturn{$interp.lookupVariable(word.name)$}\;
  }
  \uElseIf{$word.kind = \textsc{ArrayRef}$}{
    $index \leftarrow \textsc{EvalWord}(interp, word.index)$\;
    \KwReturn{$interp.lookupArrayElement(word.name, index)$}\;
  }
  \uElseIf{$word.kind = \textsc{Script}$}{
    \KwReturn{\textsc{EvalScript}($interp$, \textsc{ParseScript}(\textsc{Tokenize}($word.script$)))}\;
  }
  \Else{
    $output \leftarrow$ empty string\;
    \ForEach{$part \in word.parts$}{
      $output \leftarrow output + \textsc{EvalWord}(interp, part).asString()$\;
    }
    \KwReturn{$output$}\;
  }
}
\end{algorithm}

\section{UVec 跨设备设计}
UVec 是一个 vector-like container，但它不只拥有 CPU storage。普通 \texttt{std::vector<T>} 只维护 CPU memory；UVec 同时维护 CPU buffer 和 optional CUDA buffers，并记录哪一个 device 上的数据是 valid and up-to-date。

当 CPU code 请求最新数据，而最新副本在 CUDA 上时，UVec 执行 CUDA-to-CPU copy；当 CUDA kernel 请求最新数据，而最新副本在 CPU 上时，UVec 执行 CPU-to-CUDA copy；当某个 device 请求 mutable pointer 时，其他副本会被标记为 invalid，因为该 device 可能写入新值。图 \ref{fig:uvec-sync} 展示同步机制，图 \ref{fig:uvec-state} 展示状态机视角。

\begin{figure}[H]
  \centering
  \includegraphics[width=0.88\textwidth]{uvec-sync.png}
  \caption{UVec 跨设备同步机制（Cross-device synchronization in UVec）}
  \label{fig:uvec-sync}
\end{figure}

\begin{figure}[H]
  \centering
  \includegraphics[width=0.75\textwidth]{uvec-state.png}
  \caption{UVec 有效性状态机（UVec validity-state machine）}
  \label{fig:uvec-state}
\end{figure}

最重要的两个 pointer interfaces 是 \texttt{as\_uptr(device)} 和 \texttt{as\_mut\_uptr(device)}。前者返回 read-only pointer，并保证目标 device 上的数据已经 valid；后者返回 mutable pointer，并将目标 device 标记为 newest writer。这两个接口让 CPU algorithm 和 CUDA kernel 可以共享同一个数据抽象，而不需要脚本用户手动判断数据在哪里。

\section{形式化 UVec 一致性模型}
UVec 可以被描述为定义在 device set 上的小型 consistency model。令
\[
  D = \{\mathrm{CPU}, \mathrm{CUDA}_0, \mathrm{CUDA}_1, \ldots\}
\]
表示可能的存储设备集合。长度为 \(n\) 的 UVec instance 表示为
\[
  U = (n, B, V, L),
\]
其中 \(B_d\) 是 device \(d \in D\) 上的 physical buffer，\(V_d \in \{0,1\}\) 表示 \(B_d\) 是否包含 valid copy，\(L \subseteq D\) 表示可能包含 latest value 的 device set。当前实现通常表现为 single latest writer，但集合表示允许 CPU 和 CUDA 在一次 read synchronization 后同时 valid。

两个核心操作满足以下规则：

\begin{description}
  \item[Read rule:] \texttt{as\_uptr(d)} 只有在保证 \(V_d = 1\) 后才能返回 device \(d\) 上的 pointer。如果 \(V_d = 0\)，则从某个满足 \(V_s = 1\) 的 source device \(s\) 复制数据，并设置 \(V_d \leftarrow 1\)。Read operation 不会使其他 valid copies 失效。
  \item[Write rule:] \texttt{as\_mut\_uptr(d)} 返回 device \(d\) 上的 mutable pointer。返回前必须保证 \(B_d\) 存在且包含 latest value；随后设置 \(V_d \leftarrow 1\)，并对所有 \(s \ne d\) 设置 \(V_s \leftarrow 0\)，因为通过 mutable pointer 的后续写入可能改变 logical vector。
\end{description}

由此得到两个有用的不变式：
\[
  \exists d \in D,\; V_d = 1
\]
在每次成功初始化后的 UVec operation 之后成立；并且 UVec 返回的每个 read pointer 在返回时都指向 logical vector 的 valid copy。这个模型刻意保持简单：它足以支撑同步式 CPU/CUDA command execution，同时也为未来的 multiple CUDA devices、asynchronous streams 和 copy-on-write policies 留出空间。

\section{Geometry Command Flow}
\texttt{geom} command 的执行流程如图 \ref{fig:geom-flow} 所示。以 \texttt{geom nearest\_point} 为例，interpreter 首先求值 script arguments；command 层解析 UVec handles 和 device options；如果目标是 CPU，就把输入同步到 CPU 并调用 C++ algorithm；如果目标是 CUDA，就把输入同步到 CUDA，创建 output UVec，启动 CUDA kernel，并返回 output handle。

\begin{figure}[H]
  \centering
  \includegraphics[width=0.90\textwidth]{geom-flow.png}
  \caption{geom 命令执行流程（Execution flow of a geometry command）}
  \label{fig:geom-flow}
\end{figure}

Handle 机制非常关键。大规模点集不应该在每条命令之间反复转换为 script strings。更合理的方式是：script 创建一次 UVec，把 handle 传给多个 geometry commands，只在最终需要展示或统计时再 read back。这个思路也是后面 performance evaluation 中 speedup 能出现的前提。

\section{并发设计}
FTCL 提供 thread tasks 和 thread channels。Thread task 可以异步运行一段 script，并在之后被 await；thread channel 是由 mutex 和 condition variable 保护的 message queue，支持 blocking receive 和 non-blocking try-receive。

这个 concurrency model 对实验很有用。它可以驱动 game-like demo，模拟 producer-consumer workload，也可以通过多个 worker scripts 分发 geometry requests。当前 CUDA concurrency experiment 也暴露出一个限制：多个 script workers 可能竞争同一个 GPU，所以未来需要 streams 或 dedicated GPU scheduler。

后续第 5 章会分别从 single-GPU contention 和 multi-GPU weak-scaling 两个角度评价这个设计。这些实验也解释了为什么实现中需要把 UVec handle manager 从 global lock 细化为 handle-level locking。

\section{小结}
本章给出了 FTCL 的系统设计：语言层负责表达，command 层负责连接，UVec 层负责跨设备数据一致性，CPU/CUDA backend 负责具体计算。这个架构的关键不在于某一个单独模块，而在于这些模块之间的边界足够清楚，因而能够支持后续测试、benchmark 和应用场景扩展。

FTCL separates control, data residency, and execution backend. UVec is the bridge that keeps script-level geometry commands independent from low-level CPU/CUDA memory management.
